
\section*{Part III. Atomic physics \hfill 10 / 15}
\subsection*{1. Start from a Coulomb-bound electron}
For a hydrogen-like ion of nuclear charge \(Z\), the nonrelativistic
energy, with reduced mass \(\mu\), is
\[
\boxed{E_n=-\frac{\mu Z^2e^4}
{2(4\pi\epsilon_0)^2\hbar^2n^2}
\simeq-13.6\,{\rm eV}\,
\frac{\mu}{m_e}\frac{Z^2}{n^2}.}
\]
The eigenstates separate into radial and angular factors:
\[
\psi_{nlm}(\mathbf r)=R_{nl}(r)Y_{lm}(\theta,\phi),\quad
\int_0^\infty |R_{nl}(r)|^2r^2\,dr=1.
\]
With \(u_{nl}=rR_{nl}\), the radial Schr\"odinger equation
takes a one-dimensional form,
\[
-\frac{\hbar^2}{2\mu}\frac{d^2u}{dr^2}+
\left[V(r)+\frac{\hbar^2l(l+1)}{2\mu r^2}\right]u=Eu.
\]
The angular-momentum barrier suppresses high-\(l\) wavefunctions
near the nucleus. Hydrogenic orbitals have \(l=0,\ldots,n-1\),
\(m_l=-l,\ldots,l\), and orbital parity \((-1)^l\).
A transition photon has frequency
\(\nu=(E_i-E_f)/h\) and wavelength \(hc/(E_i-E_f)\).

\subsection*{2. Many-electron atoms: the central-field approximation}
For more than one electron, electron--electron repulsion breaks the
simple \(n\)-only degeneracy:
\[
H=\sum_i\left[\frac{p_i^2}{2m_e}
-\frac{Ze^2}{4\pi\epsilon_0r_i}\right]
+\sum_{i<j}\frac{e^2}{4\pi\epsilon_0r_{ij}}.
\]
An effective central potential replaces part of the many-body
interaction by screening. Hartree--Fock improves this by
self-consistently including direct and exchange effects in a
Slater determinant. The Aufbau \(n+l\) ordering is a useful
rule of thumb, not an exact theorem; shell energies and order
change with ionisation and configuration.

\subsection*{3. Nuclear size and relativistic effects}
A finite nuclear charge distribution changes the potential
at small \(r\), especially affecting orbitals that penetrate
the nucleus. In first-order perturbation theory,
\[
\boxed{\Delta E^{(1)}=\bra\psi\Delta V\ket\psi
=\int_0^\infty |R_{nl}(r)|^2\Delta V(r)\,r^2\,dr.}
\]
The sign and magnitude depend on the defined
\(\Delta V=V_{\rm finite}-V_{\rm point}\).
For heavy atoms, spin--orbit and other relativistic corrections
grow important; the Dirac equation incorporates the electron's
spin and relativistic kinematics from the outset.
\begin{methodbox}
\textbf{Method:} specify the central potential and boundary conditions;
choose \(R\) or \(u=rR\) consistently; solve energies/states;
evaluate corrections using normalised wavefunctions and check units.
\end{methodbox}
\src{Week 8 slides pp.~2--7, 26--36; Week 8 atomic transcripts; nuclear-size assignment.}
